\documentclass[10pt,a4paper]{amsart}
\usepackage[margin=19mm]{geometry}
\usepackage{amsmath,amssymb,mathtools}
\usepackage[scr=rsfs,cal=euler]{mathalfa}
\usepackage{xfrac}
\usepackage[unicode,hidelinks,bookmarks=false]{hyperref}
\newcommand{\ie}{i.e.\ }
\newcommand{\cf}{cf.\ }
\newcommand{\resp}{respectively\ }
\newcommand{\Subsection}{\subsection}
\providecommand{\nobreakdash}{\nobreak\hbox{-}}
\setlength{\emergencystretch}{2em}
\multlinegap0pt
\usepackage{amssymb}
\usepackage[shortlabels]{enumitem}
\usepackage{tikz}
\usepackage{xcolor}
\newtheorem{lemma}{Lemma}
\title{Law of the Iterated Logarithm for $p$-Walks on $\mathbb{Z}$}
\author{Robin Kaiser}
\date{Source: arXiv:2606.19131v1 (CC BY 4.0)}
\begin{document}
\maketitle

\noindent\emph{Source and licence.} Law of the Iterated Logarithm for $p$-Walks on $\mathbb{Z}$. Version arXiv:2606.19131v1 (CC BY 4.0), distributed by arXiv under the Creative Commons Attribution 4.0 International licence, http://creativecommons.org/licenses/by/4.0/, as recorded in the arXiv licence metadata for this exact version and shown on the abstract page https://arxiv.org/abs/2606.19131v1. Full licence notice: \texttt{licenses/arxiv-2606.19131-CC-BY-4.0.txt}.\par\smallskip

\noindent\textit{Editorial scope.} Counted unit: the article's lemma lem:variation-calculus (source0.tex lines 272--275). The article's complete original proof of it is delivered with it (source0.tex lines 276--290, one \texttt{proof} environment of 15 lines). No mathematics is written, repaired or re-derived: the statement and the proof are the article's own lines, in the article's own notation, and no retained line is substituted or re-flowed.  Retained uncounted as the source-local context the proof needs: lines 177--179.  Nothing was omitted from any retained span, so the package carries no omission record.  No retained line contains a citation, so the wrapper carries no bibliography and no bibliography entry.  No retained line contains a figure environment, an image-inclusion command or drawing code, so no image is delivered and the package declares no asset dependency.  Editorial formatting only, in addition: an AMS \texttt{amsart} preamble that restates the article's own declarations and macros used by the retained lines, namely the environments \texttt{lemma} on separate counters, together with the standard packages those lines need.  The retained environments print in this document as Lemma 1 in order of appearance, whereas the article prints the same items with the numbering of its own full document.  The complete native source of the article as distributed in the arXiv e-print for this exact version is bundled unchanged as \texttt{sources/203119/source0.tex}.
\begin{lemma}\label{lem:absolute-cont-preserved}
    The perturbation mapping $\Phi_{\alpha,\beta}$ for $\alpha,\beta<1$ preserves absolute continuity.
\end{lemma}

\begin{lemma}\label{lem:variation-calculus}
    It holds that
    $$\sup_{g\in\Phi_{\alpha,\beta}(\mathcal{K})}g(1)=\frac{1}{1-\alpha},\hspace{1cm}\inf_{g\in\Phi_{\alpha,\beta}(\mathcal{K})}g(1)=-\frac{1}{1-\beta}.$$
\end{lemma}

\begin{proof}
    We will only show how to determine the value of the supremum. The infimum works analogously. For ease of notation, let us write  $\mathcal{K}'=\Phi_{\alpha,\beta}(\mathcal{K}).$

    Let now $g\in\mathcal{K}'$, then consider the running supremum $\overline{g}$. Notice that it trivially holds that $\overline{g}(1)\geq g(1)$, thus our first goal is to show $\overline{g}\in\mathcal{K}'.$ We thus need to find a funciton $\tilde{f}\in\mathcal{K}$ such that $\overline{g}=\Phi_{\alpha,\beta}(\tilde{f})$. We define $\tilde{f}$ as the exact inverse mapping of $\overline{g}$, that is
    $\tilde{f}=(1-\alpha)\overline{g}.$
    In the proof of Lemma \ref{lem:absolute-cont-preserved} we have already seen that $\overline{g}$ must also be absolutely continuous, thus we obtain the absolute continuity of $\tilde{f}$. Since the derivative of $\tilde{f}$ agrees with the derivative of $f$ on the set $\{g=\overline{g}\}$ and is $0$ everywhere else, this implies $\tilde{f}\in\mathcal{K}.$ Since $\tilde{f}$ was defined as the exact inverse of $\overline{g}$, this implies $\overline{g}\in\mathcal{K}'$. Thus we can center our view only on the monotone functions in $\mathcal{K}'$.

    Let now $g\in\mathcal{K}'$ be a monotone increasing function. Thus it holds $\overline{g}=g$ and $\underline{g}=0$. Let us choose the pre-image $f\in\mathcal{K}$, then we have
    $f=(1-\alpha)g.$
    Thus, $g$ fulfills the following energy constraint
    $$\int_0^1(g'(t))^2dt\leq \frac{1}{(1-\alpha)^2}\int_0^1(f'(t))^2dt\leq \frac{1}{(1-\alpha)^2}.$$
    From Cauchy-Schwartz inequality and the absolute continuity of $g$, it thus follows
    $$g(1)=\int_0^1g'(t)dt\leq\left(\int_0^1(g'(t))^2dt\right)^{1/2}\leq\frac{1}{1-\alpha}.$$
    Notice that this upper bound can be achieved by the function $g(t)=\frac{1}{1-\alpha}t$, and thus the claim follows.
\end{proof}

\end{document}
